% Folder 93, batch 09 (archivists' pages 161-180).
% Pass: Opus 5.5 (claude-opus-5-5), 2026-10-10 — first pass, unchecked against the pages by a human.
%
% Contents: pp. 161-166, filtration of the relative De Rham complex and the
% Gauss-Manin connection (after Katz-Oda); p. 167 is a cover sheet in his hand
% ("Opération de Cartier"), used as a heading and given no \page; pp. 168-180,
% the Cartier operation (p-bases, Cartier isomorphism, logarithmic differentials,
% trace, duality, Verschiebung). The run continues past p. 180.
\documentclass[11pt,a4paper]{article}
\input{../preamble/grothendieck}

\folder{93}
\batch{9}
\pages{161}{180}
\foldertitle{Connexions de Gauss-Manin et équations de Picard-Fuchs. Opération de Cartier : copies de tapuscrit annoté (s.d.), tiré à part (1981), notes manuscrites (s.d.).}
\dating{[1972]-1981}
\watermark{Édition de démonstration}

\begin{document}

\section{Filtration du complexe de De Rham d'un schéma relatif lisse, et connexion de Gauss-Manin (d'après Katz-Oda)}
\note{titre de l'auteur en tête de la p.~161, numéroté « 1) » et souligné ; le mot « lisse » est ajouté au-dessus de la ligne, et un mot est biffé avant « (d'après Katz-Oda) ».}

\page{161}
1) Filtration du complexe de De Rham d'un schéma relatif\add{, lisse} et connexion de Gauss-Manin \struck{\ill{}} (d'après Katz-Oda).

\begin{tikzcd}
X \arrow[d, "f"] \\
S \arrow[d, "g"] \\
T
\end{tikzcd}

$f$ lisse [ou seulement diff.\ lisse, pour rendre exacte la suite (1)]. On a donc une suite exacte
\[
(1) \qquad \text{\struck{$0 \to$}}\; 0 \to f^{*}(\Omega^{1}_{S/T}) \to \Omega^{1}_{X/T} \to \Omega^{1}_{X/S} \to 0
\]
d'où une filtration $\mathrm{Fil}_{X/S}$ de l'objet gradué $\Omega^{\bullet}_{X/T} = \bigwedge^{\bullet} \Omega^{1}_{X/T}$, \uncertain{avec} gradué associé
\[
(2) \qquad \mathrm{Gr}^{p}_{X/S}(\Omega^{\bullet}_{X/T}) \simeq f^{*}(\Omega^{p}_{S/T}) \otimes \Omega^{\bullet}_{X/S}[-p] .
\]
On constate que la filtration respecte les opérateurs \uncertain{différentiels} de $\Omega^{\bullet}_{X/T}$, et qu'il \uncertain{en est de} même \ill{} \marginal{compte tenu que, grâce à la \uncertain{structure} différentielle de $\Omega^{\bullet}_{X/T}$ et à $f^{*}(\Omega^{p}_{S/T})$ rel.\ à $X/S$, $f^{*}(\Omega^{p}_{S/T}) \otimes \Omega^{\bullet}_{X/S}$ est \ill{} bien muni d'une \ill{} différentielle} de l'isomorphisme (2), \ill{} d'ailleurs, comme $f$ est plat, $f^{*}(\Omega^{p}_{S/T}) \simeq \mathbf{L}f^{*}(\Omega^{p}_{S/T})$, et comme $\Omega^{\bullet}_{X/S}[-p]$ est plat, le $\otimes$ dans (2) \uncertain{est} \uncertain{aussi} \uncertain{le} $\overset{\mathbf{L}}{\otimes}$, donc on a
\[
(2') \qquad \mathrm{Gr}^{p}_{X/S}(\Omega^{\bullet}_{X/T}) \simeq \mathbf{L}f^{*}(\Omega^{p}_{S/T}) \overset{\mathbf{L}}{\otimes} \Omega^{\bullet}_{X/S}[-p] ,
\]
\uncertain{dans} le \uncertain{même} \uncertain{sens} où \ill{} \uncertain{dans} \ill{}.

La filtration sur $\Omega^{\bullet}_{X/T}$ définit \marginal{$\mathbf{R}f_{*}(\Omega^{\bullet}_{X/T})$} \uncertain{donc} \ill{} \ill{} \ill{} $\Omega^{\bullet}_{X/T}$ comme objet d'une catégorie dérivée filtrée, (filtration \ill{} notée $\mathrm{Fil}^{\bullet}_{X/S}$, d'où $\mathrm{Gr}^{\bullet}_{X/S}$). On

\page{162}
aura alors
\[
\mathrm{Gr}^{p}_{X/S}\, \mathbf{R}f_{*}(\Omega^{\bullet}_{X/T}) \simeq \mathbf{R}f_{*}(\mathrm{Gr}^{p}_{X/S}(\Omega^{\bullet}_{X/T})) .
\]
Tenant compte de (2'), et supposant $f$ \uncertain{quasi-compact} \uncertain{et} \uncertain{quasi-séparé}, il résulte \add{alors} de la formule de projection [dans un cadre pas très catholique, --- il faudra faire une petite vérification]
\[
(3) \qquad \mathrm{Gr}^{p}_{X/S}\, \mathbf{R}f_{*}(\Omega^{\bullet}_{X/T}) \simeq \Omega^{p}_{S/T} \overset{\mathbf{L}}{\otimes} \mathbf{R}f_{*}(\Omega^{\bullet}_{X/S})[-p] .
\]
On a \struck{en particulier les opérateurs} une suite \uncertain{spectrale}, \uncertain{d'où} \ill{}
\[
(4) \qquad d^{(p)}_{1} : \mathrm{Gr}^{p}_{X/S}\, \mathbf{R}f_{*}(\Omega^{\bullet}_{X/T}) \to \mathrm{Gr}^{p+1}_{X/S}\, \mathbf{R}f_{*}(\Omega^{\bullet}_{X/T})[1] ,
\]
\uncertain{i.e.} \uncertain{donc}, compte tenu de l'expression explicite (3) :
\[
(5) \qquad d^{(p)}_{1} : \Omega^{p}_{S/T} \overset{\mathbf{L}}{\otimes} \mathbf{R}f_{*}(\Omega^{\bullet}_{X/S})[-p] \to \Omega^{p+1}_{S/T} \overset{\mathbf{L}}{\otimes} \mathbf{R}f_{*}(\Omega^{\bullet}_{X/S})[-p] ,
\]
\uncertain{d'où}, en translatant par $p$ :
\[
(6) \qquad \text{\struck{$\ill{}$}}\; d^{(p)}_{1}[p] : \Omega^{p}_{S/T} \overset{\mathbf{L}}{\otimes} \mathbf{R}f_{*}(\Omega^{\bullet}_{X/S}) \to \Omega^{p+1}_{S/T} \overset{\mathbf{L}}{\otimes} \mathbf{R}f_{*}(\Omega^{\bullet}_{X/S}) .
\]
On a bien \uncertain{sûr}
\[
(7) \qquad d^{(p+1)}_{1}\, d^{(p)}_{1} = 0 .
\]
Faisant $p = 0$ dans (6), on trouve
\[
(6_{0}) \qquad d^{(0)}_{1} : \mathbf{R}f_{*}(\Omega^{\bullet}_{X/S}) \to \Omega^{1}_{S/T} \otimes \mathbf{R}f_{*}(\Omega^{\bullet}_{X/S}) ,
\]
i.e.\ une connexion sur $\mathbf{R}f_{*}(\Omega^{\bullet}_{X/S})$. On l'appelle la \emph{connexion de Gauss-Manin}. \struck{[La formule (7), \uncertain{prouve} que c'est une]} (Il semble bien que c'est la \uncertain{même} que celle que j'ai définie dans \ill{} [\ill{}]!).
\marginal{Attention, il faut prouver une compatibilité de $d_{1}$ avec la diff.\ de $\Omega^{\bullet}_{S/T}$, pour expliciter comme dit (7)!}
La formule (7) montre qu'\emph{elle est intégrable}.

\page{163}
\struck{Lorsque les} Bien sûr, les différentielles (6) définissent
\[
(8) \qquad d^{p,q}_{1}[p] : \mathcal{H}^{q}(\Omega^{p}_{S/T} \otimes \mathbf{R}f_{*}(\Omega^{\bullet}_{X/S})) \to \mathcal{H}^{q}(\Omega^{p+1}_{S/T} \otimes \mathbf{R}f_{*}(\Omega^{\bullet}_{X/S})) ,
\]
qui, pour $\Omega^{\bullet}_{S/T}$ plat sur $S$, p.\ ex.\ $S$ formellement lisse sur \struck{$T$}\add{$T$}, \struck{\ill{}} s'écrit ainsi
\[
(8') \qquad d^{p,q}_{1}[p] : \Omega^{p}_{S/T} \otimes \mathbf{R}^{q}f_{*}(\Omega^{\bullet}_{X/S}) \to \Omega^{p+1}_{S/T} \otimes \mathbf{R}^{q}f_{*}(\Omega^{\bullet}_{X/S}) ,
\]
qui pour $p = 0$ définit une connexion intégrable sur les $\mathbf{R}^{q}f_{*}(\Omega^{\bullet}_{X/S})$, dite connexion de \emph{Gauss-Manin}.

Posons pour un moment, pour simplifier,
\[
(9) \qquad K^{\bullet} = \mathbf{R}f_{*}(\Omega^{\bullet}_{X/T}) ,
\]
qu'on considère \uncertain{toujours} comme complexe filtré comme il a été dit. On a \uncertain{donc}
\[
(10) \qquad \mathbf{R}\Gamma_{S}(K^{\bullet}) \simeq \mathbf{R}\Gamma_{X}(\Omega^{\bullet}_{X/T}) \simeq \mathbf{R}H_{DR}(X/T) .
\]
La filtration de $K^{\bullet}$ permet, d'après un sorite général, d'écrire une suite spectrale
\[
(11) \qquad \mathbf{R}\Gamma_{S}(K^{\bullet}) \Longleftarrow E^{p,q}_{2} = \mathbb{H}^{p}(S, \alpha \mapsto \mathcal{H}^{\alpha+q}(\mathrm{Gr}^{\alpha}(K^{\bullet}))) ,
\]
\note{le second membre est lu tel qu'écrit : « $\mathbb{H}^{p}(S, \alpha \mapsto \mathcal{H}^{\alpha+q}(\mathrm{Gr}^{\alpha}(K^{\bullet})))$ », le complexe en $\alpha$ ayant pour différentielle le $d_{1}$.}
les différentielles définissant le complexe $\alpha \mapsto \ldots$ étant les différentielles évidentes $d_{1}$ provenant de la filtration. Tenant compte de (3) \add{et (10)}, on trouve donc la suite spectrale

\page{164}
\[
(12) \qquad H^{*}_{DR}(X/T) \Longleftarrow E^{p,q}_{2} = \mathbb{H}^{p}(S, \alpha \mapsto \mathcal{H}^{q}(\Omega^{\alpha}_{S/T} \overset{\mathbf{L}}{\otimes} \mathbf{R}f_{*}(\Omega^{\bullet}_{X/S}))) ,
\]
où l'opérateur diff.\ pour le complexe $\alpha \mapsto \ldots$ s'exprime \uncertain{donc} \uncertain{directement} \uncertain{à} \uncertain{l'aide} \uncertain{de} la connexion de Gauss-Manin de $\mathbf{R}f_{*}(\Omega^{\bullet}_{X/S})$. Dans le cas favorable $\Omega^{1}_{S/T}$ plat, (il suffit $\Omega^{1}_{S/T}$ plat, i.e.\ \uncertain{Dolbeault} \uncertain{grand}) \struck{\ill{}} cette suite spectrale \uncertain{devient}
\[
(12') \qquad H^{*}_{DR}(X/T) \Longleftarrow E^{p,q}_{2} = \mathbb{H}^{p}(S, \alpha \mapsto \Omega^{\alpha}_{S/T} \otimes \mathbf{R}^{q}f_{*}(\Omega^{\bullet}_{X/S}))
\]
\[
= H^{p}_{DR}(S/T, \mathbf{R}^{q}f_{*}(\Omega^{\bullet}_{X/S})) ,
\]
l'expression $H^{p}_{DR}$ à \uncertain{droite} \uncertain{ayant} un sens, grâce à la connexion de Gauss-Manin de $\mathbf{R}^{q}f_{*}(\Omega^{\bullet}_{X/S})$.
Comme de juste, les suites spectrales (12) \struck{et (12')} (dans (12')) s'appelle\supplied{nt}, \struck{\ill{}} \emph{suite spectrale de Leray filtrée}, en cohomologie de De Rham.
\note{« suite spectrale de Leray filtrée » : lecture incertaine du dernier mot.}

\struck{\ill{}} \emph{Application}. Soit \uncertain{une} \uncertain{classe}
\[
(13) \qquad \xi \in H^{n}_{DR}(X/T)
\]
telle que son image dans $E^{0,n}_{2} = \mathbb{H}^{0}(H^{0}(S, \mathbf{R}f_{*}(\Omega^{\bullet}_{X/S})) \to H^{0}(S, \mathcal{H}^{n}(\Omega^{1}_{X/S} \overset{\mathbf{L}}{\otimes} \mathbf{R}f_{*}(\Omega^{\bullet}_{X/S}))))$ \struck{\ill{}} \uncertain{par}
\[
(14) \qquad H^{n}_{DR}(X/T) \to H^{0}(S, \mathbf{R}^{n}f_{*}(\Omega^{\bullet}_{X/S}))
\]
est \emph{nulle} [NB cette image est en tout cas toujours \emph{horizontale}, i.e.\ son image dans $\Gamma(S, \mathcal{H}^{n}(\Omega^{1}_{X/S} \overset{\mathbf{L}}{\otimes} \mathbf{R}f_{*}(\Omega^{\bullet}_{X/S})))$ est nulle], i.e.\ son image dans
\note{la formule de $E^{0,n}_{2}$ en marge droite est surchargée et en partie biffée ; on la donne telle qu'elle se laisse lire, sans garantie. Sous $H^{0}(S, \mathbf{R}^{n}f_{*}(\Omega^{\bullet}_{X/S}))$ l'auteur écrit $\mathcal{H}^{n}_{DR}(X/S)$. Les indices « $\Omega^{1}_{X/S}$ » du NB sont lus tels qu'écrits ; l'argument attend $\Omega^{1}_{S/T}$.}

\page{165}
Alors grâce à la suite spectrale de Leray, $\xi$ définit un élément dans $E^{1,n-1}_{\infty} \subset E^{1,n-1}_{2}$, donc un
\[
(15) \qquad \eta \in \mathbb{H}^{1}(S, \alpha \mapsto \mathcal{H}^{n-1}(\Omega^{\alpha}_{S/T} \overset{\mathbf{L}}{\otimes} \mathbf{R}f_{*}(\Omega^{\bullet}_{X/S}))) .
\]
Supposons en particulier les $\Omega^{1}_{S/T}$ plats, \struck{de sorte} \uncertain{donc} on obtient
\[
(16) \qquad \eta \in H^{1}_{DR}(S, \mathbf{R}^{n-1}f_{*}(\Omega^{\bullet}_{X/S})) \qquad (\Omega^{1}_{S/T} \text{ plat}) .
\]
Nous avons vu [\ill{}] qu'un tel objet s'identifie à une classe d'isomorphie de torseurs $P$ sous $\mathbf{R}^{n-1}f_{*}(\Omega^{\bullet}_{X/S})$ munis d'une \uncertain{homomorphisme}
\[
(17) \qquad P \xrightarrow{d_{p}} \Omega^{1}_{S/T} \otimes \mathbf{R}^{n-1}f_{*}(\Omega^{\bullet}_{X/S})
\]
satisfaisant
\[
(18) \qquad \begin{cases} d_{p}(x+\xi) = d_{p}(x) + d^{0}(\xi) \\ \operatorname{Im} d_{p} \subset Z^{1}(\Omega^{\bullet}_{S/T} \otimes \mathbf{R}^{n-1}f_{*}(\Omega^{\bullet}_{X/S})) . \end{cases}
\]
\marginal{$x \in P(\cdot)$, $\xi \in \mathbf{R}^{n-1}f_{*}(\Omega^{\bullet}_{X/S})(\cdot)$ ; $d^{0}$ dérivée de conn.\ de G.M.}

[On sait aussi, \uncertain{sans} \uncertain{l'hypothèse} \uncertain{des} \uncertain{platitudes} sur $\Omega^{1}_{S/T}$, que si $S$ est de car.\ nulle, la donnée de $d_{p}$ \uncertain{équivaut} \uncertain{à} \uncertain{la} \uncertain{donnée} \uncertain{de}
\[
(19) \qquad P \xrightarrow{\delta^{\infty}_{p}} P^{\infty}_{S/T}(\mathbf{R}^{n-1}f_{*}(\Omega^{\bullet}_{X/S}))^{+}
\]
satisfaisant la condition
\[
(20) \qquad \delta^{\infty}_{p}(x+\xi) = \delta^{\infty}_{p}(x) + d^{+}_{S/T}(\xi)
\]
\marginal{$x \in P(\cdot)$, $\xi \in \mathbf{R}^{n-1}f_{*}(\Omega^{\bullet}_{X/S})(\cdot)$ : $d^{+}_{S/T} = d^{\infty}_{S/T} - \delta^{\infty}_{S/T}$, $\delta^{\infty}_{S/T}$ défini par la stratification de $\mathbf{R}^{n-1}f_{*}(\Omega^{\bullet}_{X/S})$ déduite de sa connexion de G.M.}
\struck{\ill{}}
\note{le crochet ouvert devant « On sait aussi » n'est pas refermé sur la page ; l'exposant écrit « $+$ » à la fin de (19) est suivi d'un signe barré qu'on n'a pas su lire.}

\page{166}
\emph{Variantes}. a) On peut prendre sur \struck{$X$} un module de coeff.\ qu.\ coh.\ $E$, avec connexion intégrable rel.\ à $T$, d'où $\Omega^{\bullet}_{X/T} \otimes E$ avec structure de complexe. Tout ce qui précède s'applique, (ce qui \uncertain{ne} fait que la suite exacte (1) splitte localement, donne \struck{suite exacte en tensorisant les filtrations} de $\otimes$ par une filtration compatible sur $\Omega^{\bullet}_{S/T} \otimes E$, d'où
\[
\mathrm{Gr}^{p}_{X/S}(\Omega^{\bullet}_{X/T} \otimes E) \simeq \mathrm{Gr}^{p}_{X/S}(\Omega^{\bullet}_{X/T}) \otimes E
\]
au point de vue additif. Pour pouvoir appliquer la formule de projection, et obtenir, en analogie avec (3)
\[
(21) \qquad \mathrm{Gr}^{p}_{X/S}(\mathbf{R}f_{*}(\Omega^{\bullet}_{X/T} \otimes E)) \simeq \Omega^{p}_{S/T} \overset{\mathbf{L}}{\otimes} \mathbf{R}f_{*}(\Omega^{\bullet}_{X/S} \otimes E)[-p] ,
\]
il semble qu'il \struck{faille} \uncertain{suffise} \emph{$E$ plat sur $X$} \uncertain{ou} \emph{$\Omega^{1}_{S/T}$ plat sur $S$}. Sous cette réserve, on obtient des connexions intégrables sur les \struck{\ill{}} $\mathbf{R}f_{*}(\Omega^{\bullet}_{X/S} \otimes E)$ rel.\ à $T$, et une suite spectrale qui, \uncertain{dans} le 2\textsuperscript{e} cas ($\Omega^{1}_{S/T}$ plat) prend la forme simple
\[
(22) \qquad H^{*}_{DR}(X/T, E) \Longleftarrow E^{p,q}_{2} = H^{p}_{DR}(S/T, \mathcal{H}^{q}_{DR}(X/S, E)) .
\]

b) On peut localiser la suite spectrale obtenue sur $T$, (plus généralement, relativiser : un morphisme quelconque $S \to R$). On peut donner \uncertain{des} \uncertain{compatibilités} \uncertain{de} la suite spectrale obtenue sur $T$, \struck{\ill{}} \add{formules} avec des connexions intégrables relativement à une base \uncertain{au-dessus} de $T$.

\section{Opération de Cartier}
\note{titre de l'auteur, seul en haut d'un feuillet par ailleurs blanc (p.~167), qui sert de chemise à ce qui suit.}

\page{168}
\emph{L'opération de Cartier} (NB Tous les anneaux sont des $\mathbf{F}_{p}$-alg.)

① Cohomologie de De Rham d'une algèbre admettant une $p$-base.

a) Soit $B/A$, donc, posant $\Lambda = A[B^{p}] \subset B$, on a
\[
\begin{cases} \Omega^{1}_{B/A} \xrightarrow{\sim} \Omega^{1}_{B/\Lambda} & \text{d'où } \Omega^{*}_{B/A} \xrightarrow{\sim} \Omega^{*}_{B/\Lambda} \\ \text{donc } H^{*}_{DR}(B/A) \simeq H^{*}_{DR}(B/\Lambda) & \text{(comme $\Lambda$-algèbres)} \end{cases}
\]

b) Supposons que $B$ admette une $p$-base sur $A$ formée d'\emph{un} élément $x$, i.e., \struck{\ill{} \ill{}} posant $x^{p} = a \in \Lambda$, on a un isom.\ de $\Lambda$-algèbres
\[
B \simeq \Lambda[X]/(X^{p} - a) .
\]
Alors on a :

base de $\Omega^{0}_{B/\Lambda}$ sur $\Lambda$ : $1, x, x^{2}, \ldots, x^{p-2}, x^{p-1}$

base de $\Omega^{1}_{B/\Lambda}$ sur $\Lambda$ : $dx, x\,dx, x^{2}dx, \ldots, x^{p-1}dx$

d'où base de $H^{*}_{DR}(B/\Lambda) \simeq H^{*}_{DR}(B/A)$ : $\dot{1}$, $\overline{x^{p-1}dx}$.

\drawing{entre les deux lignes de bases, des flèches obliques relient $x^{k}$ à $x^{k-1}dx$ (la différentielle, avec un coefficient noté sous la troisième flèche) ; $x^{p-1}dx$ est seul à ne pas être atteint, et la flèche partant de $1$ aboutit à un point (zéro).}

c) Supposons qu'on ait une $p$-base $(x_{i})_{i \in I}$ de $B$ sur $A$. Alors \struck{\ill{}}
\[
\Omega^{1}_{B/A} \simeq \Omega^{1}_{B/\Lambda} \quad \text{admet base } (dx_{i})_{i \in I}
\]
d'où
\[
\Omega^{*}_{B/\Lambda} \simeq \bigotimes_{I,\Lambda} \Omega^{*}_{B_{i}/\Lambda} \qquad \text{(tous de $\Lambda$-alg.\ différentielles)}
\]
où
\[
B_{i} \simeq \Lambda[x_{i}] \simeq \Lambda[X_{i}]/(X_{i}^{p} - a_{i}) \qquad (a_{i} = x_{i}^{p} \in \Lambda) .
\]
Par suite, Künneth nous donne
\[
H^{*}_{DR}(B/A) \simeq H^{*}_{DR}(B/\Lambda) \simeq \bigotimes_{I,\Lambda} H^{*}_{DR}(B_{i}/\Lambda)
\]

\page{169}
② L'homomorphisme
\[
\psi : \Omega^{*}_{B/A} \to Z^{*}_{B/A} \qquad (\text{cycles de } \Omega^{*}_{B/A})
\]
Il y a un homom.\ unique $\psi$, satisfaisant les conditions

a) $\psi$ est un hom.\ d'anneaux

b) $\psi_{0}(x) = \underbrace{x^{p}}_{\psi_{0}(x)}$ si $x \in B \simeq \Omega^{0}_{B/A}$

c) $\psi(d_{B/A}x) = \underbrace{x^{p-1}d_{B/A}x}_{\partial x}$ si $x \in B$.

Pour vérifier ceci, il suffit de vérifier que l'on a
\[
\partial(xy) = \psi_{0}(x)\,\partial y + \psi_{0}(y)\,\partial x
\]
i.e.
\[
(xy)^{p-1}d(xy) = x^{p}y^{p-1}dy + y^{p}x^{p-1}dx
\]
ce qui est trivial.

Considérons $Z^{*}_{B/A}$ comme une algèbre sur $\Lambda = A[B^{p}]$, et formons
$\widetilde{B} = B \otimes_{A} \widetilde{A}$ ($\widetilde{A} = (A, F_{A})$, $F_{A} : A \to A$ défini par $x \mapsto x^{p}$),

\begin{tikzcd}[column sep=small]
B \arrow[r] \arrow[rr, bend left=30, "F_{A}"] & \widetilde{B} = B \otimes_{A} \widetilde{A} \arrow[r] \arrow[d] & B \\
A \arrow[u] \arrow[r] & \widetilde{A} = (A, F_{A}) \arrow[ur, no head] &
\end{tikzcd}

\note{diagramme redessiné d'après le croquis : la flèche verticale de gauche monte de $A$ vers $B$, celle du milieu descend de $\widetilde{B}$ vers $\widetilde{A}$, comme sur la page ; le trait de $\widetilde{A}$ vers $B$ est sans pointe. L'étiquette de l'arc courbe se lit « $F_{A}$ » ; on attendrait $F_{B}$.}

\uncertain{de} \uncertain{sorte} que l'image dans $B$ est précisément $\Lambda = A[B^{p}]$. Alors, utilisant l'homom.\ canonique $B \to \widetilde{B}$, $x \mapsto x \otimes_{A} 1_{\widetilde{A}}$, \uncertain{dont} le composé avec $\widetilde{B} \to B$ est $F_{B} : x \mapsto x^{p}$, et considérant $\Omega^{*}_{B}$ \ill{} comme une $B$-algèbre, l'homom.\ $\psi$ \uncertain{se} \ill{} \ill{} \note{le coin inférieur droit du feuillet manque ; la fin des lignes suivantes est perdue.}
\ill{} l'hom.\ $B \to \widetilde{B}$ sur les scalaires, \ill{}

\struck{un hom.\ de} \uncertain{avec} un homom.\ de $\Lambda$-\emph{algèbres}
\[
\widetilde{\psi} : \Omega^{*}_{B/A} \otimes_{B} \widetilde{B} \simeq \Omega^{*}_{\widetilde{B}/\widetilde{A}} \to Z^{*} \ldots
\]
\struck{qui est donc \ill{}} Nous serons intéressés surtout par l'homomorphisme composé
\[
Z^{*}_{B/A} \to H^{*}_{DR}(B/A) \ldots
\]
savoir
\[
\boxed{\varphi : \Omega^{*}_{\widetilde{B}/\widetilde{A}} \to H^{*}_{DR}(B/A)}
\]

\page{170}
③ L'opération de Cartier.

\emph{Théorème 1}. Supposons que $B/A$ admette une $p$-base. Alors l'homom.\ précédent $\varphi$ \add{induit} \struck{est} un \emph{isomorphisme} $\Omega^{*}_{\widetilde{B}/\widetilde{A}} \otimes_{\widetilde{B}} \Lambda \to H^{*}_{DR}(B/\Lambda)$ [NB $\Lambda \simeq \widetilde{B}/\text{idéal}$] donc est un isom.\ si $\widetilde{B} \xrightarrow{\sim} \Lambda$ i.e.\ $\widetilde{B} \to B$ est injectif.

En effet, si $(x_{i})_{i \in I}$ est une $p$-base, utilisant les \uncertain{notations} du n\textsuperscript{o}~1, on trouve un isom.\ de $B$-algèbres
\[
\Omega^{*}_{B/A} \simeq \bigotimes_{i \in I} \Omega^{*}_{B_{i}/\Lambda} \otimes_{B_{i}} B
\]
d'où \marginal{isom.\ de $B$-alg.}
\[
\Omega^{*}_{\widetilde{B}/\widetilde{A}} \simeq \bigotimes_{i \in I} (\Omega^{*}_{\widetilde{B}_{i}/\widetilde{\Lambda}} \otimes_{\widetilde{B}_{i}} \widetilde{B})
\]
et on connaît de même \uncertain{la} \uncertain{structure} \uncertain{de} $H^{*}_{DR}(B/A)$ :
\[
H^{*}_{DR}(B/A) \simeq \bigotimes_{i \in I} H^{*}_{DR}(B_{i}/\Lambda) .
\]
L'homom.\ $\varphi$ envisagé s'identifie alors au produit tensoriel des homomorphismes analogues \struck{correspondants}
\[
\varphi_{i} : \Omega^{*}_{\widetilde{B}_{i}/\widetilde{\Lambda}} \otimes_{\widetilde{B}_{i}} \widetilde{B}_{i} \to H^{*}_{DR}(B_{i}/\Lambda)
\]
\[
\text{base } \widetilde{1},\ d\widetilde{x}_{i} \qquad\qquad \text{base } \dot{1},\ \overline{x_{i}^{p-1}dx_{i}}
\]
qui envoient base sur base, donc sont des isom.\ cqfd pour le théorème.

\begin{tikzcd}[column sep=small, row sep=small]
B \arrow[r] & \widetilde{B} \arrow[r] & B \\
B_{i} \arrow[u] \arrow[r, no head] & \widetilde{B}_{i} \arrow[u, no head] \arrow[r, no head] & B_{i} \arrow[u] \\
\Lambda \arrow[u] \arrow[r, no head] & \widetilde{\Lambda} \arrow[u, no head] \arrow[r, no head] & \Lambda \arrow[u] \\
A \arrow[u, no head] \arrow[r, no head] & \widetilde{A} \arrow[u, no head] &
\end{tikzcd}

\note{diagramme en marge gauche, recopié tel que dessiné : seules les flèches de la première ligne et les petites flèches verticales ont des pointes nettes ; les autres traits sont rendus sans pointe.}

\emph{Corollaire 1}. On a \struck{\ill{}}
\[
Z^{*}(B/A) \simeq B^{*}(B/A) \oplus [\Omega^{*}_{\widetilde{B}/\widetilde{A}} \otimes_{\widetilde{B}} \Lambda]
\]

\emph{Définition}. L'isom.\ inverse de $\varphi$ est noté
\[
\boxed{\widetilde{C} : H^{*}_{DR}(B/A) \xrightarrow{\sim} \Omega^{*}_{\widetilde{B}/\widetilde{A}} \otimes_{\widetilde{B}} \Lambda \ (= \Omega^{*}_{B/A} \otimes_{B} \Lambda)}
\]
On dira aussi souvent $C$ l'homom.
\[
\widetilde{C} : Z^{*}(B/A) \to \Omega^{*}_{\widetilde{B}/\widetilde{A}}
\]
donnant les précédents par passage aux quotients, [dans ce \uncertain{cas} il est \emph{surjectif}, \emph{et son noyau est formé des différentielles exactes}].

\page{171}
\emph{Corollaire 2}. Il y a une application unique
\[
\widetilde{C} : Z^{*}(B/A) \to \Omega^{*}_{\widetilde{B}/\widetilde{A}}
\]
satisfaisant aux conditions suivantes

a) $\widetilde{C}(d\omega) = 0$ pour tout $\omega \in Z^{*}(B/A)$

b) $\widetilde{C}$ est un hom.\ de \add{$\Lambda$-\emph{algèbres}} \struck{anneaux}
\[
\begin{cases} \widetilde{C}(\omega + \omega') = C(\omega) + C(\omega') \\ \widetilde{C}(\omega\omega') = C(\omega)\,C(\omega') \\ \widetilde{C}(\lambda\omega) = \lambda\,C(\omega) \quad \text{si } \lambda \in \Lambda \\ \widetilde{C}(1) = 1 \end{cases}
\]

c) $\widetilde{C}(x^{p-1}d_{B/A}x) = d_{\widetilde{B}/\widetilde{A}}x \otimes_{\widetilde{B}} 1_{\Lambda}$ \qquad d) $\widetilde{C}\left(\dfrac{dx}{x}\right) = \dfrac{dx}{x} \otimes_{B} 1_{\Lambda} = \dfrac{d\widetilde{x}}{\widetilde{x}} \otimes_{\widetilde{B}} 1_{\Lambda}$

\note{dans c), un symbole est noyé d'encre devant $x^{p-1}$. La dernière condition (le cas $\widetilde{C}(1) = 1$) est entourée, avec la précédente, d'un trait reliant les deux.}

L'existence résulte de la définition. L'unicité résulte du corollaire 1.

\emph{Corollaire 3}. Si $x \in B^{*}$, on a d).

En effet, on a $\dfrac{dx}{x} = x^{-p}(x^{p-1}dx)$, et comme $x^{-p} \in \Lambda$, on trouve
\[
\widetilde{C}\left(\frac{dx}{x}\right) = x^{-p}\,\widetilde{C}(x^{p-1}dx) = \text{\struck{$\ill{}$}}\; F_{A}(x^{-1})\,C(x^{p-1}dx)
\]
\[
= (\widetilde{x}^{-1}d\widetilde{x}) \otimes_{\widetilde{B}} \Lambda
\]

\emph{Cas particuliers}

a) Supposons que $\operatorname{Im}(A \to B) \subset B^{p}$, \struck{i.e.\ $A \to \operatorname{Im}(A \to B)$. Alors} \uncertain{par} \uncertain{exemple} $A$ parfait. Alors $\Omega^{*}_{B/A} \simeq \Omega^{*}_{B/\mathbf{F}_{p}} \overset{\text{déf}}{=} \Omega^{*}_{B}$, et \uncertain{de} \uncertain{même} pour $H^{*}_{DR}(B/A)$, de sorte que \uncertain{on} \uncertain{peut} supposer que $A = \mathbf{F}_{p}$. Mais dans ce dernier cas, on a $A \xrightarrow{\sim} \widetilde{A}$, donc $B \simeq \widetilde{B}$, et si $B \xrightarrow{F_{B}} B$ est injectif i.e.\ $\widetilde{B} \xrightarrow{\sim} \Lambda$ i.e.\ $B$ \emph{réduit}, alors \uncertain{le} \uncertain{th.} \uncertain{1} \uncertain{devient}
\[
\Omega^{*}_{B} \xrightarrow[\sim]{\varphi} H^{*}_{DR}(B) .
\]
\uncertain{D'où} l'opération de Cartier \uncertain{peut} être considérée \uncertain{alors} comme
\[
\widetilde{C} : Z^{*}(B) \to \Omega^{*}_{B}
\]
On \uncertain{fait} \uncertain{attention} que la \add{dernière} formule s'écrit ici

b') \qquad $\widetilde{C}(f^{p}\omega) = f\,\widetilde{C}(\omega)$

et c) et d) deviennent

\page{172}
\[
\begin{cases} C(x^{p-1}dx) = dx \\ C(dx/x) = dx/x \end{cases}
\]
C'est le cas envisagé dans le \uncertain{texte} initial de Cartier.

④ \uncertain{Fonctorialités}, Globalisation.
\note{« Fonctorialités » est relié par un trait à « le cas envisagé » de la ligne au-dessus ; les deux mots semblent avoir été intervertis ou rapprochés après coup.}

Soit $X \to S$ en car.\ $p$, d'où

\begin{tikzcd}
X \arrow[d] & X^{(p/S)} \arrow[l] \arrow[d, "{f = f_{X/S}}"] & X \arrow[l] \arrow[ll, bend right=30, "F_{X}"'] \arrow[dl, no head] \\
S & S \arrow[l, "F_{S}"] &
\end{tikzcd}

\note{dans le carré de gauche, l'auteur écrit « cart. » (souligné) : carré cartésien. Le trait de $X$ (à droite) vers le $S$ du milieu est sans pointe sur la page.}

avec $f = f_{X/S}$ radiciel \add{entier} surjectif].
\struck{[Supposons que $f_{*}(\underline{O}_{X})$ soit localement \ill{} \ill{} $p$-base ; i.e.\ $\operatorname{Im}(\underline{O}_{X^{(p/S)}} \to f_{*}(\underline{O}_{X}))$, alors on trouve un isom.\ \ill{}.]}
On définit donc un hom.\ canonique de faisceaux d'algèbres
\[
\psi : \Omega^{*}_{X^{(p/S)}/S} \to f_{*}(Z^{*}(\Omega^{*}_{X/S}))
\]
d'où un hom.\ canonique
\[
\psi : \Omega^{*}_{X^{(p/S)}/S} \to f_{*}(\mathcal{H}^{*}(\Omega^{*}_{X/S})) \simeq \mathcal{H}^{*}(f_{*}(\Omega^{*}_{X/S})) ,
\]
\note{dans le dernier membre, un symbole biffé devant $(\Omega^{*}_{X/S})$.}
Si $\underline{\Lambda} = \operatorname{Im}(\underline{O}_{X^{(p/S)}} \to f_{*}(\underline{O}_{X}))$ [qui est donc de la forme $\underline{O}_{X^{(p/S)}}/\mathcal{J}$, $\mathcal{J}$ un nilidéal, donc $\mathcal{J} = 0$ si $X^{(p/S)}$ est réduit, p.\ ex.\ $S$ réduit et $X$ \struck{\ill{}} \add{loc.}\ de présentation finie sur $S$], on trouve donc
\[
\widetilde{\psi} : \Omega^{*}_{X^{(p/S)}/S} \otimes_{\underline{O}_{X^{(p/S)}}} \underline{\Lambda} \to \text{\struck{$f_{*}$}}\; \mathcal{H}^{*}(f_{*}(\Omega^{*}_{X/S}))
\]
Si $f_{*}(\underline{O}_{X})$ admet loc.\ une $p$-base sur $\underline{O}_{X^{(p/S)}}$ (ou, ce qui revient au \uncertain{même}, sur $\underline{\Lambda}$), alors l'hom.\ précédent est un isom., et \uncertain{on} \uncertain{peut} définir, \struck{\ill{}} l'isom.\ inverse
\[
\widetilde{C} : \mathcal{H}^{*}(f_{*}(\Omega^{*}_{X/S})) \to \Omega^{*}_{X^{(p/S)}/S} \otimes_{\underline{O}_{X^{(p/S)}}} \underline{\Lambda}
\]
d'où un hom.
\[
\widetilde{C} : Z^{*}(f_{*}(\Omega^{*}_{X/S})) \to \Omega^{*}_{X^{(p/S)}/S} \otimes_{\underline{O}_{X^{(p/S)}}} \underline{\Lambda}
\]
\marginal{définir $\widetilde{C}$ chaque fois que $\psi$ est un isom.\ (sans qu'il y ait nécessairement des $p$-bases …)}

\page{173}
La formation de $\psi$, chaque fois $\widetilde{C}$ sous ses conditions d'existence, satisfait à une propriété de fonctorialité évidente pour un carré commutatif

\begin{tikzcd}
X \arrow[d] & X' \arrow[l] \arrow[d] \\
S & S' \arrow[l]
\end{tikzcd}

donnant naissance à un diagramme commutatif.

\begin{tikzcd}[column sep=small, row sep=small]
 & X \arrow[dl] \arrow[ddl] & & X' \arrow[ll] \arrow[dl] \arrow[ddl] \\
X^{(p/S)} \arrow[d] & & X'^{(p/S')} \arrow[ll] \arrow[d] & \\
S & & S' \arrow[ll] &
\end{tikzcd}

\note{le cube est redessiné d'après le croquis ; on a rendu les flèches telles qu'elles paraissent, de $X$ vers $X^{(p/S)}$ et vers $S$, et de même pour $X'$.}


⑤ \struck{Relation avec les} \emph{Connexions sur les Modules inversibles, et différentielles logarithmiques}.

Soit $\underline{L}$ module inversible sur $X$, \struck{\ill{}} \add{muni d'une connexion $\Delta_{0}$ rel.\ à $S$.} Alors les connexions $\Delta$ sur $\underline{L}$, \struck{qui} \add{qui correspondent} (rel.\ à $S$), \struck{s'identifient aux} homom.\ linéaires $u : \underline{L} \to \underline{L} \otimes \Omega^{1}_{X/S}$ \add{par $u \mapsto \Delta = \Delta_{0} + u$}, s'identifient aux 1-formes différentielles \struck{\ill{}} $\omega$ sur $X/S$.
Si $\Delta_{0}$ est à courbure nulle, \uncertain{la} connexion sur $\underline{L}$ définie par $\omega$ \uncertain{est} à courbure nulle, \uncertain{ssi} $d\omega = 0$.

\emph{Proposition}. Pour que la connexion sur $\underline{L}$ définie par $\omega$ soit \struck{\ill{}} \add{à courbure nulle}, il faut et il suffit que l'on ait $d\omega = 0$.
\struck{Pour qu'elle soit} Supposons que $X/X^{(p/S)}$ soit \add{loc.}\ \add{à $\Delta$, $\Delta_{0}$ à courbure nulle, avec $\Delta_{0}$ comp.\ aux puiss.\ $p$-ièmes} \struck{\ill{}} [Pour que la connexion \struck{\ill{}} admette des $p$-bases, et que \ill{} \ill{} \uncertain{compatible} avec \uncertain{les} puissances $p$-ièmes envisagée \uncertain{comme} [i.e.\ $\Theta^{p}_{\xi}(x) = \Theta_{\xi^{p}}(x)$ pour tout champ de vecteurs $\xi$ \uncertain{sur} \ill{}, et \ill{} section \uncertain{de} $\underline{L}$ \uncertain{sur} celui-ci)
\note{la fin de la page (à partir de « Supposons que $X/X^{(p/S)}$ ») est très surchargée : insertions interlinéaires croisées et crochets non refermés ; on donne ce qui se lit, l'ordre des insertions restant incertain. L'énoncé se poursuit p.~174.}

\page{174}
il faut et il suffit que l'on ait $\widetilde{C}\omega = \widetilde{\omega} \otimes_{\underline{O}_{X^{(p/S)}}} 1_{\underline{\Lambda}}$.

La première assertion est standard, et ne dépend pas de l'hyp.\ de car.\ $p$. [Si $\underline{L}$ \uncertain{était} \uncertain{pas} de rang 1, mais \uncertain{seulement} localement \uncertain{libre}, alors $\omega$ est une forme diff.\ à valeurs dans $\underline{\mathrm{End}}(E)$, et la relation à écrire est $d\omega + [\omega, \omega] = 0$.]

Pour la deuxième, la condition étant locale, on peut supposer \uncertain{que} $X$ est affine. On a donc \uncertain{que} \ill{} \ill{} \uncertain{admet} \uncertain{une} $p$-base $(x_{i})$, d'où \uncertain{des} $D_{i} = \dfrac{\partial}{\partial x_{i}}$.
\struck{$\ill{}\ \rho(D_{i}) = \rho_{0}(D_{i}) + \ill{}$}
\uncertain{D'où} Écrivons la condition que $\Delta$, \struck{\ill{}} \add{qui donne} $\xi \mapsto \Theta_{\xi}$, \struck{$\rho(D_{i})^{p}$} est compatible aux puissances $p$-ièmes. On trouve \struck{$\Theta_{\ill{}}$} \ill{} des conditions nécessaires
\[
\Theta_{D_{i}}^{p} = 0 \qquad (i \in I)
\]
\uncertain{car} $D_{i}^{p} = 0$. Je dis qu'elles sont suffisantes. \uncertain{Cela} \struck{\ill{}} \uncertain{se} voit aisément, en \uncertain{remarquant} que l'algèbre des \struck{\ill{} \ill{} \ill{} \ill{}} \uncertain{opérateurs} \struck{$\sum a_{i}D_{i}$,} \uncertain{effet}, \struck{\ill{} \ill{}} \ill{} différentiels engendrés par les $D_{i} = \dfrac{\partial}{\partial x_{i}}$ et les mult.\ \struck{les $a_{i} \in B$.} scalaires est engendrée par \add{les quantités} $D_{i}$, $x_{i}$, avec \add{les relat.}
\[
\begin{cases} [D_{i}, D_{j}] = 0 \\ [x_{i}, x_{j}] = 0 \\ [D_{i}, x_{j}] = \delta_{ij} \\ x_{i}^{p} = a_{i} \\ D_{i}^{p} = 0 \end{cases}
\]
Il faut donc, compte tenu de
\[
\Delta_{i} = \Theta_{D_{i}} = \Theta^{0}_{D_{i}} + \omega(D_{i}).\mathrm{id}
\]
(où $\Delta^{0}_{i}$ correspond à $\Theta^{0}_{D_{i}}$), écrivant $\omega = \sum c_{i}\,dx_{i}$, d'où $c_{i} = \omega(D_{i})$, exprimer la condition
\[
(\Delta^{0}_{i} + c_{i})^{p} = 0 .
\]
Mais la formule de \uncertain{Jacobson} donne

\marginal{Notes en diagonale dans la marge gauche, en partie illisibles : « \uncertain{maladroit} \uncertain{des} \uncertain{choix} \uncertain{des} \uncertain{bases} » (souligné) ; « Dire : la \uncertain{relation} de \uncertain{courbure} $\Theta^{p}$ \uncertain{définissant} la \uncertain{compatibilité} \uncertain{aux} puissances $p$-ièmes \uncertain{comme} une $\Theta^{p}$-\uncertain{courbure} \ill{} $\Omega^{1}_{X/S} \otimes \underline{\mathrm{End}}(E)$ \uncertain{des} anneaux \uncertain{envisagés} […] \ill{} \uncertain{nulle} \uncertain{sur} \uncertain{les} $\Theta$-\uncertain{courbure} \uncertain{donnée} » ; « \ill{} [\uncertain{alors} \ill{}] \uncertain{égale}, \uncertain{plus}, $\widetilde{C}\omega = F^{*}_{X/S}(\widetilde{\omega} - C\omega)$ … $\omega \oplus p - F_{X/S}$ \uncertain{nulle} \ill{} \ill{} \uncertain{donc} $\widetilde{\omega} = \widetilde{C}\omega$ » ; « \uncertain{Indépendante} \uncertain{de} l'hyp.\ que $L$ \uncertain{de} rang 1 ».}

\page{175}
(cf.\ p.\ ex.\ Sém.\ Chevalley 58/59, exp.~6 p.~06)
\[
(\Delta^{0}_{i} + c_{i})^{p} = \underbrace{\Delta^{0\,p}_{i}}_{=\,0 \text{ par hyp.}} + c_{i}^{p} + D_{i}^{p-1}c_{i}
\]
donc la condition envisagée s'écrit
\[
(*) \qquad c_{i}^{p} + D_{i}^{p-1}c_{i} = 0 \qquad i \in I
\]
Or écrivons que \add{$d\omega = 0$,} \struck{$\widetilde{C}\omega = \widetilde{\omega}$, en trouvant} ($\omega = \sum c_{i}\,dx_{i}$) on trouve pour $d\omega = 0$ la condition
\[
(**) \qquad D_{i}c_{j} - D_{j}c_{i} = 0 \qquad i, j \in I
\]
et on peut alors écrire
\[
(***) \qquad \widetilde{C}\omega = \text{\struck{$\ill{}$}}\; \sum_{i} -D_{i}^{p-1}c_{i}\,(d\widetilde{x}_{i} \otimes 1_{\Lambda})
\]
\marginal{sous $D_{i}^{p-1}c_{i}$ : « est dans $\Lambda$ car égal à $\langle \widetilde{D}_{i}, C\omega \rangle 1_{\Lambda}$ ».}
\struck{ce qui signifie aussi que $\sum (c_{i}^{p} + D_{i}^{p-1}c_{i})\,x_{i}^{p-1}\,dx_{i} \sim$ cohomologue à $\sum c_{i}^{p}\,dx_{i}$, i.e.\ que $\sum c_{i}^{p}$ \ill{}}
[ceci : la formule de transposition \uncertain{incluse} dans \uncertain{la} \uncertain{formulation} \uncertain{des} §3--§4 du loc.\ cit.\ p.~05 (prop.~3)]
\note{les trois mots en diagonale de la marge (« incluse dans … formulation … des §3--§4 ») sont rattachés par des traits au crochet ; leur place exacte est incertaine.}
alors que
\[
\widetilde{\omega} = \sum c_{i}^{p}\,(d\widetilde{x}_{i} \otimes 1_{\Lambda})
\]
et on trouve donc,
\[
-\widetilde{C}\omega + \widetilde{\omega} = \sum (c_{i}^{p} + D_{i}^{p-1}c_{i})(d\widetilde{x}_{i} \otimes 1_{\Lambda})
\]
donc c'est nul ssi (*) l'est, cqfd.

Nous utiliserons maintenant le fait bien connu qu'une connexion à courbure nulle \struck{dont} \add{et} compatible à la puissance $p$-ième, sur un Module quasi-cohérent, équivaut à une \emph{donnée de descente} de $X$ à $(X^{(p/S)}, \underline{\Lambda})$. \struck{On trouve donc}

\page{176}
\struck{loc.\ sur $X^{(p/S)}$, $\underline{O}_{X}$} Prenons $\underline{L} = \underline{O}_{X}$, avec $\Delta_{0}$ la connexion triviale. Alors $\omega$ satisfaisant $d\omega = 0$, $\widetilde{C}\omega = \widetilde{\omega}$ définit donc une donnée de descente sur $\underline{L}$, qui \struck{\ill{}} provient de $\underline{L}_{0}$ sur $X^{(p/S)}_{0} = (X^{(p/S)}, \underline{\Lambda})$. Localement, on peut alors \struck{\ill{}} trouver une base de $\underline{L}_{0}$, \struck{\ill{}} \uncertain{ce} qui \uncertain{signifie} \uncertain{que} \uncertain{loc.} on a une section \uncertain{inversible} engendrant $\underline{L}$ par une section horizontale $f$, ce qui signifie que loc.\ on a ainsi les $\omega = \dfrac{df}{f}$. On trouve ainsi le

\emph{Théorème}. Pour que (\add{la 1-forme}) $\omega$ soit loc.\ une dérivée logarithmique, i.e.\ que $\omega = df/f$ ($f \in \Gamma(\underline{O}_{X}^{*})$) il faut et il suffit que l'on ait $d\omega = 0$, \add{$C\omega = \widetilde{\omega}$.} En d'autres termes, on a une \uncertain{suite} \uncertain{exacte}
\[
\underline{O}_{X}^{*} \xrightarrow{f \mapsto df/f} \underline{\Omega}^{1}_{X/S}
\]
dont l'image est le sous-faisceau de $\underline{\Omega}^{1}_{X/S}$ formé des $\omega$ telles que $d\omega = 0$, $\widetilde{C}\omega = \widetilde{\omega}$. [NB le noyau est $\underline{\Lambda}^{*}$, donc si $S$ est parfait, c'est $\underline{O}_{X}^{*\,p}$]

⑥ \emph{Cas des schémas relat.\ lisses. Relations avec la théorie des formes diff.\ de degré maximum.}

Supposons $X$ lisse sur $S$, donc $X^{(p/S)}$ lisse sur $S$. Alors $X \xrightarrow{F_{X/S}} X^{(p/S)}$ est un morphisme fini localement libre (en plus d'être radiciel et surjectif), à

\page{177}
fortiori $\underline{O}_{X^{(p/S)}} \xrightarrow{\sim} \underline{\Lambda}$ \struck{Donc On sait bien} [Il suffit \add{\uncertain{en} \uncertain{effet} \uncertain{que} $X_{s} \to X_{s}$ \uncertain{est} \uncertain{fini}, $X_{s}$ (\uncertain{étant}) \uncertain{régulier}, \uncertain{à} \uncertain{dimension} \uncertain{plus} \uncertain{source}} de vérifier fibre par fibre … ] On sait bien que $X/S$ a localement une $p$-base i.e.\ qu'il en est ainsi de $X$ sur $X^{(p/S)}$ : en effet, \uncertain{il} \uncertain{suffit} \uncertain{de} \uncertain{savoir} que $X$ est loc.\ de \uncertain{tf.}\ sur $S$, cette condition équivalant à celle que $\Omega^{1}_{X/S}$ soit localement libre, ce qui est bien le cas !

On est donc dans un cas où l'opération de Cartier est
\[
\widetilde{C} : F_{X/S\,*}(Z\Omega^{*}_{X/S}) \to \Omega^{*}_{X^{(p/S)}}
\]
Si la dim.\ relative est $n$, \uncertain{on} trouve en particulier
\[
\widetilde{C} : F_{X/S\,*}(\underbrace{\underline{\Omega}^{n}_{X/S}}_{\substack{\omega_{X/S} \\ \text{module dualisant relatif}}}) \to \underbrace{\underline{\Omega}^{n}_{X^{(p/S)}}}_{\omega_{X^{(p/S)}}}
\]

\emph{Théorème 3}. L'homomorphisme précédent est l'homomorphisme trace (ou de Gysin) $F_{X/S\,*}$ associé au morphisme $F_{X/S} : X \to X^{(p/S)}$ de schémas lisses sur $S$.

Soient donc $f_{1}, \ldots, f_{n}$ une suite de sections de $\underline{O}_{X}$ qui définissent un sous-schéma $Z$ \add{(donc loc.\ lib.\ sur $S$)} de $X$ \emph{fini} sur $S$. \uncertain{Les} $\widetilde{f}_{1}, \ldots, \widetilde{f}_{n}$ définissent \uncertain{donc} $\widetilde{Z}$ item, et \uncertain{pour} toute forme, \struck{Pour} une section $\overline{\omega}$ de $\Omega^{n}_{X^{(p/S)}}$, \struck{$\mathrm{Tr}$} $\operatorname{res}\left(\dfrac{\overline{\omega}}{\widetilde{f}_{1} \ldots \widetilde{f}_{n}}\right)$.
On peut aussi \uncertain{prendre} les images $f_{1}^{p}, \ldots, f_{n}^{p}$ dans $X$, qui définissent un sous-schéma $Z'$ de $X$ fini sur $S$, et pour les $n$-formes $\omega$ sur $X/S$, on a les formes $\operatorname{res}\left(\dfrac{\omega}{f_{1}^{p} \ldots f_{n}^{p}}\right)$. On trouve alors, comme \uncertain{formule} essentielle

\page{178}
\uncertain{revient} équivalemment d'exprimer le th.~3, la relation
\[
\boxed{\operatorname{res}\left(\frac{\widetilde{C}\omega}{\widetilde{f}_{1} \ldots \widetilde{f}_{n}}\right) = \operatorname{res}\left(\frac{\omega}{f_{1}^{p} \ldots f_{n}^{p}}\right)}
\]
Lorsque la dim.\ relative est 1, cette formule s'écrit
\[
\operatorname{res}\frac{\widetilde{C}\omega}{\widetilde{f}} = \operatorname{res}\frac{\omega}{f^{p}} \quad \text{i.e.} \quad \boxed{\operatorname{res}\widetilde{C}\overline{\omega} = \operatorname{res}\overline{\omega}}
\]
\uncertain{si} \uncertain{on} pose $\overline{\omega} = \omega/f^{p}$ ; \uncertain{le} premier membre \uncertain{est} \uncertain{aussi} $\operatorname{res}\widetilde{C}(\omega/f^{p})$.
Et lorsque de plus $S$ est parfait i.e.\ $F_{S} : S \to S$ un isom., de sorte que $C$ est défini, on peut l'écrire
\[
\text{\struck{$\left(\operatorname{res}\dfrac{C\omega}{f}\right)^{p} = \operatorname{res}\dfrac{\omega}{f^{p}}$}} \qquad \boxed{(\operatorname{res} C\overline{\omega})^{p} = \operatorname{res}\overline{\omega}}
\]
\note{sous le premier membre biffé, l'auteur relie $\operatorname{res}\dfrac{C\omega}{f}$ à $\operatorname{res} C\!\left(\dfrac{\omega}{f^{p}}\right)$ par un signe $=$.}

\emph{Remarques}. ① Il y a aussi un hom.\ de Cartier
\[
F_{i} : F_{*}(\Omega^{i}_{X/S}) \to \Omega^{i}_{X^{(p/S)}/S} ,
\]
défini grâce à l'h.\ \uncertain{trace} de Cartier \uncertain{pour} $i = n$ et la formule de projection. Mais comme
\[
F^{*} : \Omega^{i}_{X^{(p/S)}/S} \to F_{*}\Omega^{i}_{X/S}
\]
est nul pour $i \neq 0$, on conclut que $F_{i} = 0$ si $i \neq n$. En particulier, si $i \neq 0$, l'hom.\ $F_{i}$ ne coïncide \add{pas en général} avec l'opération de Cartier \add{sur les formes fermées}.

② Il y a lieu de définir aussi l'opération de Cartier \struck{\ill{}} sur les modules dualisants relatifs $\underline{\omega}_{X/S}$, \add{sans hyp.\ de lissité,} lorsque la formation de celui-ci commute aux chgt.\ de base $S \xrightarrow{F_{S}} S$ (p.\ ex.\ $S$ parfait i.e.\ $F_{S}$ un isom.), \struck{\ill{}} $X$ étant loc.\ de type fini sur $S$ loc.\ noethérien. On trouve
\[
C = F_{X/S\,*} : F_{X/S\,*}(\omega_{X/S}) \to \omega_{X^{(p/S)}/S} \simeq \widetilde{\omega}_{X/S}
\]
d'où, \uncertain{en} \uncertain{appliquant} $F^{-1}_{X/S}$ et composant avec $F^{-1}_{X/S}(\widetilde{\omega}_{X/S}) \to \widetilde{\omega}_{X/S}$, \uncertain{un} hom.\ (pas linéaire, mais $\Lambda$-linéaire)
\[
\omega_{X/S} \to \omega^{(p)}_{X/S}
\]
\note{la flèche finale porte en exposant « $(p)$ » ; on la lit $\omega_{X/S} \to \omega^{(p)}_{X/S}$, sans garantie.}

\page{179}
Appliquons le th.~3 à la dualité globale, en supposant $X$ \emph{propre et lisse} sur $S$. Soit $E$ module loc.\ libre sur $X$ (\uncertain{et} \uncertain{supposons} $S$ parfait …), $\widetilde{E}$ son image inverse sur $X^{(p/S)}$, \add{dont l'image inverse sur $X$ par $F_{X/S}$ est $E^{(p)}$.} Alors on a des homomorphismes
\[
\text{\struck{$[H^{i}(X, E) \to]\ H^{i}(X^{(p/S)}, \widetilde{E}) \to H^{i}(X, E^{(p)})$}}
\]
\uncertain{dit} « \uncertain{puissance} \uncertain{$p$-ième} »,
\[
(\varphi) \qquad \mathbf{R}f_{*}(\widetilde{E}) \xrightarrow{\ \varphi\ } \mathbf{R}f_{*}(E^{(p)})
\]
\[
\phantom{(\varphi) \qquad} \simeq \mathbf{R}f_{*}(E)^{(p)}
\]
qui n'est autre essentiellement que l'hom.\ de changement de base associé au carré cartésien

\begin{tikzcd}
X \arrow[d, "f"'] \arrow[r, "F_{X}"] & X \arrow[d] \\
S \arrow[r, "F_{S}"] & S
\end{tikzcd}

\note{les flèches horizontales du carré sont dessinées sans pointe, de gauche à droite ; on les rend dans le sens où elles sont écrites. À côté du $X$ de droite : $F_{X}^{*}(E) = E^{(p)}$.}

Considérons le \emph{transposé} de ($\varphi$), nous savons que c'est l'hom.
\[
\mathbf{R}f_{*}(\check{E}^{(p)} \otimes \Omega^{n}_{X/S}[n]) \to \mathbf{R}f_{*}(\check{\widetilde{E}} \otimes \Omega^{n}_{\widetilde{X}/S}[n])
\]
\note{au-dessus du premier membre, l'auteur écrit « $\simeq \mathbf{R}g_{*}(\mathbf{R}F_{X/S\,*}(E^{(p)} \otimes \Omega^{n}_{X/S}))$ ».}
associé à l'hom.\ trace
\[
(**) \qquad F_{X/S\,*}(\check{E}^{(p)} \otimes \Omega^{n}_{X/S}) = \check{\widetilde{E}} \otimes F_{X/S\,*}(\Omega^{n}_{X/S})
\]
\[
\phantom{(**) \qquad} \to \check{\widetilde{E}} \otimes \Omega^{n}_{X^{(p/S)}/S}
\]
\uncertain{dirons}, par tensorisation avec $\check{\widetilde{E}}[n]$ \struck{\ill{}} de l'hom.\ trace habituel. Celui-ci à son tour s'interprète comme l'hom.\ \struck{trace} de Cartier.

En particulier, on voit que le transposé de
\[
R^{n}f_{*}(\widetilde{E}) \longrightarrow R^{n}f_{*}(E^{(p)})
\]
\[
\phantom{R^{n}f_{*}(\widetilde{E})} \simeq R^{n}f_{*}(E)^{(p)}
\]
\struck{[\ill{} \uncertain{supposons} $R^{n}f_{*}(E)$ et $R^{n}f_{*}(E^{(p)})$ loc.\ libres, \ill{} \ill{}]}

\page{180}
est l'hom.
\[
f_{*}(\check{E}^{(p)} \otimes \Omega^{n}_{X/S}) \to g_{*}(\check{\widetilde{E}} \otimes \Omega^{n}_{X^{(p/S)}/S})
\]
\[
\phantom{f_{*}(\check{E}^{(p)} \otimes \Omega^{n}_{X/S}) \to{}} \simeq g_{*}(\check{E} \otimes \Omega^{n}_{X/S})^{(p)}
\]
\note{sous l'isomorphisme vertical : « sous réserve de commutation à $F_{S}$ de $f_{*}(\check{E} \otimes \Omega^{n}_{X/S})$ » ; un premier membre $f_{*}(\check{E} \otimes \Omega)$ est biffé à gauche.}
donné par l'opération de Cartier (**), ci-dessus.

Le cas le plus important est celui où $E = \underline{O}_{X}$, on trouve que l'hom.\ de puissance $p$-ième
\[
R^{i}f_{*}(\underline{O}_{X})^{(p)} \to R^{i}f_{*}(\underline{O}_{X})
\]
est transposé de l'opération de Cartier
\[
\widetilde{C} : f_{*}(\Omega^{n}_{X/S}) \to g_{*}(\Omega^{n}_{X^{(p/S)}/S}) \underset{\text{s.\ réserve}}{\simeq} f_{*}(\Omega^{n}_{X/S})^{\otimes p}
\]
\note{l'exposant final se lit « $\otimes p$ » ; on attendrait $(p)$.}
$\overline{M}$ \uncertain{énoncé} pour les $\underbrace{R^{i}f_{*}(\underline{O}_{X})^{(p)}}_{\text{sous réserve que c'est } R^{i}f_{*}(\underline{O}_{X^{(p/S)}})} \to R^{i}f_{*}(\underline{O}_{X})$, \uncertain{à} \uncertain{cela} \uncertain{près} qu'il faut prendre l'hom.\ sur
\[
R^{n-i}f_{*}(\Omega^{n}_{X/S}) \to R^{n-i}f_{*}(\Omega^{n}_{X/S})^{(p)}
\]
\note{sous ce dernier exposant : « [sous réserve] ».}
\uncertain{déduit} de $\widetilde{C}$ par \uncertain{application} de $R^{n-i}f_{*}$.

Un hom.\ intéressant est
\[
R^{1}f_{*}(\underline{O}_{X^{(p/S)}}) \longrightarrow R^{1}f_{*}(\underline{O}_{X})
\]
dans le cas où $f$ est propre, \struck{car} \struck{\ill{}} il s'identifie à l'hom.\ des \uncertain{sections} pour les \uncertain{plans} \uncertain{tangents}
\[
\underline{\mathrm{Lie}}\,\underline{\mathrm{Pic}}_{X^{(p/S)}/S} \longrightarrow \underline{\mathrm{Lie}}\,\underline{\mathrm{Pic}}_{X/S}
\]
\[
\Vert
\]
\[
(\underline{\mathrm{Lie}}\,\underline{\mathrm{Pic}}_{X/S})^{(p/S)}
\]
Nous admettrons ici que l'hom.\ \struck{\ill{}}
\[
\underline{\mathrm{Pic}}_{X^{(p/S)}/S} \simeq \underline{\mathrm{Pic}}^{(p/S)}_{X/S} \xrightarrow{\ V\ } \underline{\mathrm{Pic}}_{X/S}
\]
est le \emph{Verschiebung}, d'où il résulte que l'hom.
\marginal{Il est immédiat en tout cas que $F_{\mathrm{Pic}}V = p\,\mathrm{id}$, $VF_{\mathrm{Pic}} = p\,\mathrm{id}$ \struck{\ill{}}.}
\marginal{À vérifier : cela doit donner tout aussi des \uncertain{conditions} \uncertain{intéressantes} \uncertain{des} $p$-\uncertain{rang} \ill{}.}
\note{l'argument se poursuit au-delà de la p.~180, hors de ce lot.}

\end{document}
